%% 第 15 章：复特征值对应的实二维不变子空间
%% 本文件由 tools/md2tex.py 自动生成，请勿直接编辑。

\chapter{复特征值对应的实二维不变子空间}


\section{从复特征向量得到实不变平面}

设：

\[
z=x+iy
\]

是 $T_{\mathbb C}$ 的特征向量，且：

\[
T_{\mathbb C}(z)
=
(a+ib)z,
\qquad
b\ne0.
\]

比较实部与虚部，得到：

\[
T(x)=ax-by,
\]

\[
T(y)=bx+ay.
\]

因此：

\[
\operatorname{span}_{\mathbb R}(x,y)
\]

是 $T$ 的二维实不变子空间。

\section{对应的实矩阵块}

在实基：

\[
(x,y)
\]

下，$T$ 在该二维不变子空间上的矩阵为：

\[
\begin{pmatrix}
a&b\\
-b&a
\end{pmatrix}.
\]

这个二维块同时描述：

\begin{jitem}
\item %
$a$ 所代表的缩放；
\item %
$b$ 所代表的旋转耦合。
\end{jitem}

若 $b=0$，它退化为实特征值对应的一维结构。

\section{复 Jordan 块如何回到实空间}

对于一对共轭复 Jordan 块：

\[
J_r(a+ib),
\qquad
J_r(a-ib),
\]

回到实空间后，它们合并为一个实块矩阵：

\begin{jitem}
\item %
每个复标量 $a+ib$ 替换为：

\[
\begin{pmatrix}
a&b\\
-b&a
\end{pmatrix};
\]
\item %
原 Jordan 块超对角线上的 $1$ 替换为：

\[
I_2.
\]
\end{jitem}

这得到实 Jordan 结构。

\section{复化的意义}

复化让实算子的真实谱结构在更大的数域中显现，再将共轭成对的复结构无损翻译为实空间中的二维不变块。
